% Junior Level 3, Week 5 -- workbook chapter (read aloud by a grown-up)
\chapterweek{5}{Play Day: Blocks to Words}{Every Scratch block has a word twin in Python.}

\section{What you will learn}
\begin{ican}
  \item I can match a Scratch block to its Python word twin.
  \item I can find a bug and fix it.
  \item I can plan a pattern in parts and build it.
\end{ican}

\begin{exercise}[Blocks to Words]
Python is a coding language with words instead of blocks. Draw a line from each block to its word twin.
\begin{center}
\renewcommand{\arraystretch}{2.4}
\begin{tabular}{@{}l p{2.2cm} l@{}}
\begin{scratch}[scale=0.8]\blockmove{move \ovalnum{100} steps}\end{scratch} & $\bullet$\hfill$\bullet$ & \texttt{for i in range(4):}\\
\begin{scratch}[scale=0.8]\blockmove{turn \turnleft{} \ovalnum{90} degrees}\end{scratch} & $\bullet$\hfill$\bullet$ & \texttt{t.color("blue")}\\
\begin{scratch}[scale=0.8]\blockrepeat{repeat \ovalnum{4}}{\blockspace[0.4]}\end{scratch} & $\bullet$\hfill$\bullet$ & \texttt{t.forward(100)}\\
\begin{scratch}[scale=0.8]\blockpen{set pen color to \pencolor{kcblue}}\end{scratch} blue & $\bullet$\hfill$\bullet$ & \texttt{t.left(90)}\\
\end{tabular}
\end{center}
The \texttt{t} is the turtle that draws in Python.
\end{exercise}

\begin{exercise}[Bug-Hunt Relay]
Three projects are broken. Click the green flag. Find the bug. Fix it. Get a sticker!

\textbf{Leg 1: The Window.} It should draw four squares.
\begin{center}
\begin{scratch}[scale=0.8]
\blockinit{when \greenflag clicked}
\blockpen{erase all}
\blockpen{pen down}
\blockrepeat{repeat \ovalnum{4}}{
  \blockrepeat{repeat \ovalnum{4}}{
    \blockmove{move \ovalnum{60} steps}
    \blockmove{turn \turnright{} \ovalnum{90} degrees}
    \blockmove{turn \turnright{} \ovalnum{90} degrees}
  }
}
\end{scratch}
\end{center}

\textbf{Leg 2: The Magic Trick.} It should say \enquote{Watch this!}, then vanish.
\begin{center}
\begin{scratch}[scale=0.8]
\blockinit{when \greenflag clicked}
\blocklook{show}
\blocklook{say \ovalnum{Watch this!} for \ovalnum{2} seconds}
\end{scratch}
\qquad
\begin{scratch}[scale=0.8]
\blockinit{when \greenflag clicked}
\blocklook{hide}
\end{scratch}
\end{center}

\textbf{Leg 3: The Stuck Counter.} It should count every click on the apple.
\begin{center}
\begin{scratch}[scale=0.8]
\blockinit{when this sprite clicked}
\blockvariable{set \selectmenu{apples} to \ovalnum{0}}
\blockvariable{change \selectmenu{apples} by \ovalnum{1}}
\end{scratch}
\end{center}
\end{exercise}

\begin{exercise}[My Best Pattern]
Plan first. Write or draw your three parts.
\begin{enumerate}[leftmargin=1.4em,itemsep=2pt]
  \item Part 1, get ready: \rule{6cm}{0.4pt}
  \item Part 2, the pattern: \rule{6cm}{0.4pt}
  \item Part 3, the ending: \rule{6cm}{0.4pt}
\end{enumerate}
Build one part. Test it. Then build the next part. Your pattern needs a repeat inside a repeat. Here is one idea.
\begin{center}
\begin{scratch}[scale=0.8]
\blockinit{when \greenflag clicked}
\blockpen{erase all}
\blockpen{set pen color to \pencolor{kcpurple}}
\blockpen{pen down}
\blockrepeat{repeat \ovalnum{6}}{
  \blockrepeat{repeat \ovalnum{4}}{
    \blockmove{move \ovalnum{50} steps}
    \blockmove{turn \turnright{} \ovalnum{90} degrees}
  }
  \blockmove{turn \turnright{} \ovalnum{60} degrees}
}
\blocklook{say \ovalnum{Ta-da!} for \ovalnum{2} seconds}
\end{scratch}
\end{center}
\end{exercise}

\begin{grownup}[Show me]
While your child plays, use the Week 5 show-me checklist in the Assessments booklet. It is not a test. \enquote{Not yet} is fine.
\end{grownup}

\begin{vocab}
  \vocabword{Python}{A coding language that uses words instead of blocks.}
  \vocabword{word twin}{A Python line that does the same job as a Scratch block.}
\end{vocab}

\begin{reflection}
Color one star for each relay leg you fixed.

\begin{center}
\Huge\faStar[regular]\qquad\faStar[regular]\qquad\faStar[regular]
\end{center}

How do you feel about Play Day? Circle one.

\begin{center}
\Huge\faSmile[regular]\qquad\faMeh[regular]\qquad\faFrown[regular]
\end{center}
\end{reflection}
